\documentclass[11pt]{article}
\usepackage[T1]{fontenc}
\usepackage[utf8]{inputenc}
\usepackage{lmodern,microtype}
\usepackage[margin=1in]{geometry}
\usepackage{amsmath,amssymb,amsthm,mathtools}
\usepackage{booktabs,longtable,array,enumitem}
\usepackage{xcolor,hyperref,fancyhdr,listings}
\definecolor{inkblue}{RGB}{25,54,77}
\hypersetup{colorlinks=true,linkcolor=inkblue,citecolor=inkblue,urlcolor=inkblue,
 pdftitle={Marked Determinant Continuations and Singular-Safe Boundary Compression},
 pdfauthor={OpenAI, research assistance for the ProveIt project}}
\setlength{\parskip}{4pt}
\setlength{\parindent}{0pt}
\setlength{\emergencystretch}{2em}
\setlist{nosep,leftmargin=1.5em}
\pagestyle{fancy}
\fancyhf{}
\fancyhead[L]{\small Marked determinant continuations}
\fancyhead[R]{\small ProveIt research continuation}
\fancyfoot[C]{\thepage}
\setlength{\headheight}{14pt}
\newtheorem{theorem}{Theorem}[section]
\newtheorem{proposition}[theorem]{Proposition}
\newtheorem{lemma}[theorem]{Lemma}
\newtheorem{corollary}[theorem]{Corollary}
\theoremstyle{definition}
\newtheorem{definition}[theorem]{Definition}
\newtheorem{example}[theorem]{Example}
\theoremstyle{remark}
\newtheorem{remark}[theorem]{Remark}
\newcommand{\F}{{\mathbb F_2}}
\newcommand{\Z}{\mathbb Z}
\newcommand{\Q}{\mathbb Q}
\newcommand{\N}{\mathbb N}
\newcommand{\Kh}{\operatorname{Kh}}
\newcommand{\rk}{\operatorname{rank}}
\newcommand{\corank}{\operatorname{corank}}
\newcommand{\poly}{\operatorname{poly}}
\newcommand{\sig}{\sigma_4}
\newcommand{\nuFour}{\nu_4}
\newcommand{\norm}[1]{\left\lVert#1\right\rVert_1}
\newcommand{\code}[1]{\texttt{#1}}
\newcommand{\zero}{\mathbf 0}
\lstset{basicstyle=\ttfamily\small,breaklines=true,columns=fullflexible,
 keepspaces=true,frame=single,rulecolor=\color{black!20},aboveskip=8pt,belowskip=8pt}
\title{\color{inkblue}\bfseries Marked Determinant Continuations\\[4pt]
\Large and Singular-Safe Boundary Compression\\[7pt]
\large Recovered quantum gradings, residue-four witnesses,\\
critical square laws, and stopping-prefix complexity}
\author{Research assistance prepared for the ProveIt project}
\date{7 October 2026}

\begin{document}
\maketitle
\thispagestyle{empty}
\begin{abstract}
We develop a read-only obstruction layer for the repository's exact
Bar--Natan scanning algorithm. The layer evaluates actual completed direct
summands rather than only the original knot's determinant. A retained marked
arc makes reduced completion functorial. Homogeneous differential entries
recover the erased quantum shifts of each connected component, uniquely up
to a constant, so the existing ungraded component keys suffice. The
$\ell^1$ norm of the quantum Euler polynomial modulo $q^4-1$ then gives a
monotone lower bound on reduced Khovanov rank. Its matching-level values are
computable by link connectivity and a signed Tait determinant, without
expanding the remaining crossing cube.

For repeated determinant queries we prove a singular-safe terminal reduction.
An interior block of corank $r$ is retained as a zero block coupled to the
boundary; each nonzero query uses a matrix of order at most $2b-2$, where
$b$ is the terminal count. At the critical number $r+1$ of terminal blocks,
all nonzero quotient determinants lie in one signed rational square class.
The corank threshold is sharp. Exact Python implementations, independent
state-sum and homology oracles, and paired graph-kernel measurements accompany
the proofs.

The resulting quasi-polynomial estimate is a theorem about successful
prefixes with bounded physical representation size and frontier width.
No universal bound on those parameters, unrestricted quasi-polynomial
recognizer, or end-to-end production speedup is claimed.
\end{abstract}
\vfill
{\small\textbf{Status.} Research manuscript with written proofs and executable
checks; not peer reviewed or formally verified. The package is an integration
candidate, not a modification of the remote repository. Exact scope and
remaining obligations are stated in Sections~\ref{sec:scope} and~\ref{sec:validation}.}
\clearpage
\tableofcontents
\clearpage

\section{Research target and position relative to earlier work}
\label{sec:scope}

\subsection{The bottleneck being addressed}
The project in \path{Topology/UnknotRecognition} already has an exact scanning
backend, literal direct-sum sharing, scalar-unit cancellation, and an
Euler-based completed-component observer. Recent continuations have addressed
twist blocks, homogeneous coefficient multiplication, radical residue
profiles, and exact homological windows. Those are starting points, not
results claimed again here. The reviewed component implementation stores a
matching identifier, a relative homological degree, and every differential
coefficient in a complete key~\cite{repo,component}. Its multiplicities are
ordinary nonnegative integers, not elements of $\F$.

We address a different question:
\begin{quote}
Can one extract a stronger, inexpensive obstruction from the actual
remaining diagram and the already discovered direct summands, without
computing the suffix's homology or changing the existing ungraded sharing
keys?
\end{quote}
The answer is affirmative under precise relative-complex assumptions. The
observer combines the $q=1$ and $q=i$ specializations of the reduced quantum
Euler polynomial. The former is controlled by link connectivity; the latter
is a signed determinant. The norm is taken \emph{after summing within each
true chain summand and before summing across independent summands}. This
order prevents cancellations between independent homology groups from hiding
an obstruction.

A second question is algorithmic. Many matching completions can share a large
signed graph interior. Ordinary electrical-network elimination suggests a
boundary response matrix, but signed weights invalidate the automatic
invertibility of the interior block. We retain exactly the missing nullspace
coupling. The result is a small determinant formula, a sharp vanishing gate,
and a square-class law for critical terminal partitions.

\subsection{What is classical and what is established here}
The foundations are Khovanov's categorification~\cite{khovanov},
Bar--Natan's tangle formalism and local cancellation~\cite{bn-tangles,bn-fast},
and the classical relation between Jones specializations and signed graph
matrices; see also Boninger~\cite{boninger}. Schur complements and Kron
reduction are established tools~\cite{kron}. We do not claim these ingredients
as discoveries.

The contributions developed and proved in this continuation are the
following precise statements and their compatibility with this scanner:
\begin{enumerate}
\item Connected ungraded component keys determine quantum shift differences;
therefore they also determine the completed residue norm, despite forgetting
absolute quantum grading.
\item A mark-preserving completion theorem gives a monotone residue-four
rank obstruction, with an exact formula for the rank it fails to see.
\item A signed terminal kernel remains valid when the interior is singular.
Its corank gate is sharp, and critical partitions obey a signed square law.
\item These operations have explicit output-sensitive costs, with a
quasi-polynomial consequence for controlled successful scan prefixes.
\end{enumerate}
The mathematical arguments are supplied in full. A comprehensive novelty or
priority search has not been completed, and no theorem is presented as an
independently reviewed resolution of the unrestricted recognition problem.

\subsection{Current complexity status}
Lackenby's July 2026 preprint gives a hierarchy-based incompressibility and
unknot-recognition algorithm. Section 9 explicitly distinguishes estimates
for the number of steps from a running-time estimate for implementations
whose operations are not yet specified at that level~\cite[Section 9]{lackenby}.
Our algebraic observer is not an implementation of that hierarchy algorithm.

Kronheimer--Mrowka prove that reduced Khovanov rank one detects the
unknot~\cite{km}. For the early rejection result here, only the simpler fact
$\dim_{\F}\widetilde{\Kh}(U)=1$ and invariance are needed. A lower bound greater
than one rules out the unknot. Equality of a \emph{lower bound} with one does
not establish equality of the actual rank.

Three different meanings of ``small'' must remain separate. A small frontier
limits the size of one coefficient algebra. A small number of matching types
does not bound their multiplicities. A small observed component key does not
bound the number of physically materialized intermediate objects. Our cost
statement uses the latter quantity explicitly, including temporary expansion
before cancellation.

\subsection{Source audit and reproducibility boundary}
The observed repository head was
\begin{center}\small\texttt{8e5edcedb580523bdd4946233f7d0b42a97901d4}.\end{center}
The current component key and the Conway input were inspected at that
revision. The scanner allocation and recovered-grading diagnostic were
inspected at
\begin{center}\small\texttt{dfc762327721d9d8040eb62a36a4ab20c005b9e5}.\end{center}
Per-file blob hashes and the exact scope of each read are recorded in
\path{provenance.json}. This is a source audit, not a claim that an entire
checkout of either revision was tested. The local independent package was
executed; the supplied upstream diagnostic was not.

\section{Recovering quantum information from ungraded components}
\label{sec:grading}

\subsection{Objects, coefficients, and weights}
Fix a valid frontier with $2k$ endpoints. Matching objects have $k$ arcs.
For two objects $a,b$, let $c(a,b)$ be the number of circles in their overlay.
A coefficient is a binary linear combination of dotted basis morphisms,
indexed by subsets of these circles. If a subset $S$ contains $d$ circles,
its intrinsic weight is
\begin{equation}\label{eq:weight}
 \omega(x_S:a\longrightarrow b)=k-c(a,b)+2d.
\end{equation}
The negative of this weight is the intrinsic quantum degree of the
cobordism in the scanner's convention. A nonzero homogeneous coefficient has
one common dot count $d$ across all its monomials.

Write $h_v$ for the homological degree and $q_v$ for the quantum shift of
object $v$. A differential entry $f_{uv}$ satisfies $h_v=h_u+1$ and
\begin{equation}\label{eq:potential}
 q_v-q_u=k-c(a_u,a_v)+2d(f_{uv}).
\end{equation}
The repository's grading diagnostic uses precisely this identity, rearranged
as $2d=c-k+q_v-q_u$~\cite{grading}.

\begin{lemma}[Preservation of homogeneity]\label{lem:homogeneity}
Start from a genuine graded crossing complex and apply graded delooping and
scalar-unit cancellations, while erasing shifts only from storage. Every
remaining nonzero differential coefficient is homogeneous. A unit entry
between equal matching objects is the undotted identity.
\end{lemma}
\begin{proof}
The crossing complex and the Frobenius delooping maps are homogeneous.
Composition adds intrinsic degrees. A coefficient between two fixed shifted
objects can therefore contain only monomials of their prescribed degree.
For equal matchings $c(a,a)=k$. If a homogeneous unit has a scalar term, that
term has dot count zero and forces $q_v-q_u=0$ in
\eqref{eq:potential}. Every other monomial would have positive dot count and
thus a different degree. No such term can occur, so the unit is exactly the
identity. Cancelling this entry is a homogeneous change of complex: the
usual Schur correction is a composition of homogeneous maps of the required
degree. Induction proves the assertion throughout genuine cancellations.
\end{proof}

This lemma does not license applying the procedure to arbitrary fabricated
bit matrices. The validity of the relative complex remains a precondition.
The implementation checks coefficient homogeneity and potential consistency
but does not reprove the provenance of the entire input complex.

\subsection{The quantum-potential theorem}
\begin{theorem}[Recovery and uniqueness]\label{thm:quantum}
Let $T$ be a connected nonzero-differential component of a genuine partial
complex. Its matching objects and homogeneous coefficients determine all
quantum shifts up to one additive integer constant. One can recover them by
assigning zero to one vertex and propagating~\eqref{eq:potential} along the
undirected differential graph. Every remaining edge gives a consistency
check.

If two connected components have equal complete ungraded serialized keys,
the key's vertex correspondence is a graded chain isomorphism after a single
uniform quantum shift, in addition to the uniform homological shift already
allowed by the key.
\end{theorem}
\begin{proof}
Orient each nonzero differential edge from its source to its target and give
it the integer weight on the right-hand side of~\eqref{eq:potential}. Give
the reverse edge the negative weight. Genuine quantum shifts solve these
equations, so the sum of weights around every closed walk is zero. The sum
along a path from the chosen root is consequently independent of the path.
It is the desired potential after subtracting the root's true shift.

If $q$ and $q'$ are two solutions, their difference has zero change along
every edge. Connectedness makes that difference constant. A one-vertex
component is included: its single shift is unconstrained but immaterial up
to translation.

Equal ungraded keys preserve matching objects, nonzero edges, coefficients,
and relative homological degrees. In particular they preserve every weight
in~\eqref{eq:potential}. Pulling the second component's potential back to
the first therefore gives another solution, which differs by one constant.
With that shift the explicit vertex correspondence preserves both gradings
and every differential entry.
\end{proof}

\begin{corollary}[No new absolute grading key is required]\label{cor:key}
Any quantity on completed components invariant under uniform homological and
quantum shifts can be evaluated from the existing complete ungraded keys.
It is unnecessary to expand a shared multiplicity into its unknown absolute
quantum-shift distribution.
\end{corollary}
\begin{proof}
Theorem~\ref{thm:quantum} identifies all copies with a chosen recovered graded
representative up to precisely the shifts under which the quantity is
invariant. Each copy contributes the same value.
\end{proof}

\paragraph{Cost and implementation.}
If there are $s$ vertices and $e$ stored differential entries, propagation
uses $O(s+e)$ graph operations after edge weights have been obtained.
Homogeneity checking scans the actual coefficient support. It is not a
constant-time operation on exponentially long bit-packed values. The
observer therefore retains the scanner's coefficient-size factor in its
complexity bound. The function \code{recover\_quantum\_shifts} checks cycles
and rejects inhomogeneous coefficients; it does not silently choose the
lowest-degree term of a mixed coefficient.

\section{Marked completion and residue-four rank witnesses}
\label{sec:marked}

\subsection{Why the marked arc cannot be omitted}
Let $A=\F[x]/(x^2)$, with the usual Khovanov Frobenius grading. A marked point
on a link gives an $A$-action by adding a dot to the marked component. We use
the reduced complex normalized so that a single marked circle contributes
one generator in quantum degree zero. Equivalently, use the quotient by the
marked $x$-action with its conventional overall shift, or the marked-$x$
subcomplex with the opposite normalization shift.

At scan stage $i$, let $C_i$ be the processed relative tangle complex. Its
actual remaining diagram defines a completion operation $F_i$. Choose a
point on a small arc in the unprocessed suffix, disjoint from all processed
operations. Denote completion followed by marked reduction by
$F_i^\bullet$. A practical choice is to keep one root crossing last and mark
an incident segment close to that crossing. All proper prefixes retain it.

\begin{theorem}[Marked completion preserves relative equivalences]
\label{thm:mark}
Suppose every chain map and chain homotopy in a relative equivalence of
processed tangle complexes is extended by the identity near the retained
marked arc. Marked completion preserves that chain-homotopy equivalence and
preserves finite direct sums.
\end{theorem}
\begin{proof}
After completion, every map commutes with the dot action at the fixed marked
arc: the map is the identity on a neighborhood of that arc, and the
Frobenius construction is compatible with its marked $A$-action. The maps
and homotopies are therefore $A$-linear. Suppose $f,g$ and homotopies $H,K$
satisfy
\[
 gf-1=dH+Hd,\qquad fg-1=dK+Kd.
\]
Passing to the marked quotient applies an additive functor to each identity.
The same equations still hold on the reduced complexes. No assertion that
the quotient functor is exact, or that its coefficient module is flat, is
required: these are preserved explicit homotopy equations. Completion and
the quotient both commute with finite direct sums. Uniform grading
normalization does not alter the argument.
\end{proof}

\begin{remark}[A forbidden shortcut]\label{rem:divide}
It is incorrect to take arbitrary summands of a fully delooped, unmarked
closed complex and divide each summand's Euler polynomial by
$q+q^{-1}$. A single surviving unmarked generator need not be a free marked
$A$-summand. The total complex may admit a marked reduction without the
chosen unmarked splitting respecting it. Keeping the mark in the suffix is
what makes the direct-sum statement above legitimate. The provided observer
requires a nonempty marked suffix and declines a final unmarked closure.
\end{remark}

\subsection{An integral four-coordinate shadow}
For a finite bigraded cochain complex $X$ with quantum-degree-preserving
differential, define
\[
 \chi_X(q)=\sum_{h,j}(-1)^h\dim_\F X^{h,j}\,q^j.
\]
Reduction modulo $q^4-1$ gives a unique vector in $\Z^4$:
\begin{equation}\label{eq:sigma}
 \sig(X)_r=\sum_{h,\ j\equiv r\ (4)}(-1)^h\dim_\F X^{h,j},
 \qquad \nuFour(X)=\norm{\sig(X)}.
\end{equation}
These are \emph{integer} alternating sums even though the differential is
computed over $\F$.

\begin{lemma}[Shadow norm properties]\label{lem:shadow}
The vector $\sig$ is invariant under graded chain-homotopy equivalence and
additive under direct sums. Its norm satisfies
\[
 0\le\nuFour(X)\le\dim_\F H(X).
\]
A uniform homological shift changes the sign of the vector when odd; a
uniform quantum shift permutes its coordinates. Both preserve $\nuFour$.
\end{lemma}
\begin{proof}
For each quantum degree, the alternating dimension of a finite complex
agrees with that of its homology, by cancellation of boundary dimensions in
adjacent homological degrees. Summing by quantum residue proves invariance.
Additivity is immediate. In each residue class, the absolute value of the
alternating sum is no greater than the sum of the homology dimensions;
summing proves the inequality. The shift statements follow directly
from~\eqref{eq:sigma}.
\end{proof}

\begin{theorem}[Marked completed-component obstruction]\label{thm:bound}
Assume the original input is a validated classical knot $K$, and a genuine
relative decomposition at stage $i$ has the form
\[
 C_i\simeq\bigoplus_a\ \bigoplus_{h,j}
 T_a[h]\{j\}^{\oplus p_{a,h,j}},\qquad p_{a,h,j}\in\N.
\]
Set $m_a=\sum_{h,j}p_{a,h,j}$. Then
\begin{equation}\label{eq:bound}
 B_4(i)=\sum_a m_a\nuFour(F_i^\bullet T_a)
 \ \le\ \dim_\F\widetilde{\Kh}(K;\F).
\end{equation}
In particular $B_4(i)>1$ certifies that $K$ is nontrivial. The norm of each
completed template can be evaluated using its recovered quantum potentials;
absolute quantum shifts of shared copies are unnecessary.
\end{theorem}
\begin{proof}
Theorem~\ref{thm:mark} transports the displayed relative decomposition to a
direct sum of reduced completed complexes for $K$, up to the irrelevant
global diagram normalization. Homology of this finite direct sum is the
direct sum of its homologies. Each shifted copy has the same total homology
dimension and, by Lemma~\ref{lem:shadow}, the same shadow norm as its
representative. Apply the norm bound separately to every copy and sum.
The normalized reduced homology of the unknot has dimension one, proving
rejection when the lower bound is larger. Finally
Corollary~\ref{cor:key} applies because the norm is shift-invariant.
\end{proof}

\begin{proposition}[Monotonicity under refinement]\label{prop:monotone}
Suppose one scan step extends each current representative additively,
performs genuine relative cancellations, and splits the result into direct
summands. The resulting $B_4$ cannot decrease while the mark is retained.
Exact sharing of equal keys does not change the value.
\end{proposition}
\begin{proof}
For a parent $T$, completing before or after that scan step gives homotopy
equivalent reduced complexes. If the extended parent splits into children
$S_j$ with shifts, its shadow equals the sum of the shifted child shadows.
By the triangle inequality, its norm is at most the sum of their norms.
Multiply by the parent's nonnegative multiplicity and sum over parents.
Merging equal child keys merely adds the numbers of identical contributions.
\end{proof}

An arbitrary change of direct-sum decomposition need not be a refinement,
so the proposition must not be asserted for every unrelated decomposition
chosen by another algorithm. The maintained component engine's parent-local
extension and splitting is the relevant situation.

\subsection{Comparison with simpler observers}
For $v=(v_0,v_1,v_2,v_3)$, define
\[
 e(v)=\left|\sum_r v_r\right|,
 \qquad d(v)=|v_0-v_2|+|v_1-v_3|.
\]
The first is the absolute Euler characteristic. The second is the
$\ell^1$ norm of the Gaussian integer obtained by evaluating at $i$.
Both are bounded by $\norm{v}$. Hence $B_4$ dominates the corresponding
\emph{reduced marked} completed-component Euler bound and the determinant
character bound. This does not identify it numerically with a particular
unreduced observer in the production pipeline; normalizations must be
compared at the same marked level.

If a component's Euler polynomial has only one parity of quantum exponents,
its four-vector has only two possible nonzero coordinates. For their sum
$a$ and difference $d$, its norm is
\begin{equation}\label{eq:max}
 \left|\frac{a+d}{2}\right|+
 \left|\frac{a-d}{2}\right|=\max(|a|,|d|).
\end{equation}
In general the four-vector formulation is safer than imposing a parity
promise on an arbitrary formal component.

\begin{proposition}[Exact invisible-rank formula]\label{prop:defect}
Let $E_r$ and $O_r$ be the total homology dimensions in quantum residue $r$
and even and odd homological degrees, respectively. Then
\begin{equation}\label{eq:defect}
 \dim H(X)-\nuFour(X)=2\sum_{r=0}^3\min(E_r,O_r).
\end{equation}
Consequently the shadow is exact if and only if, in every quantum residue,
homology occurs in at most one homological parity.
\end{proposition}
\begin{proof}
The residue coordinate is $E_r-O_r$. Subtract $|E_r-O_r|$ from $E_r+O_r$;
the difference is $2\min(E_r,O_r)$. Sum over residues. All summands are
nonnegative, giving the equality criterion.
\end{proof}

\subsection{Strictness, with a realization boundary}
\begin{example}[Independent summands can reveal hidden rank]\label{ex:strict}
Let $I$ be a two-ended interval object. Consider the dotted two-term complex
\[
 T=\bigl(I\xrightarrow{x}I\{2\}\bigr)
\]
in homological degrees zero and one. Complete with a marked crossingless
arc. On marked reduction the dotted map is zero, so its Euler polynomial is
$1-q^2$ and its reduced rank is two. The admissible relative complex
\[
 I\ \oplus\ T\ \oplus\ T[1]\{4\}
\]
has global residue-four norm one but split residue-four bound five. Its
split ordinary Euler bound is one because each two-term block has Euler
characteristic zero. The dotted blocks are connected in the differential
graph before marked reduction.

This is a strict example in the relative algebraic category. No claim is
made that this particular complex is the invariant of a knot tangle. In
particular it is not a construction of a nontrivial knot with trivial Jones
polynomial. Its purpose is to prove that the proposed norm and the position
at which it is applied can matter strictly.
\end{example}

\paragraph{Safe saturation.}
For a rejection threshold of two on reduced rank, nonnegative multiplicities
and final nonnegative contributions may be capped at two. If the existing
backend supplies multiplicities capped at three, those are still safe lower
bounds. Signed coordinates must be accumulated exactly until the norm is
taken. Replacing $M-(M-1)$ by independently capped terms can turn a true
value of one into zero or otherwise change a subtraction; saturation is not
a homomorphism on signed integer arithmetic.

There is a further distinction for a saturated scanner. If every multiplicity
is capped at a fixed integer $c\ge2$, then
\[
 \min\!\left(2,\sum_a\min(c,m_a)\nuFour(F_i^\bullet T_a)\right)
 =\min(2,B_4(i)).
\]
Indeed, any changed nonzero product already reaches two; unchanged products
are exact. Thus the threshold-capped value remains monotone. The
\emph{uncapped sum made from capped multiplicities} need not be monotone:
merging two contributions of multiplicity three can replace their sum six
by a shared capped multiplicity three. It is still a valid lower bound.
The observer reports which monotonicity statement applies and includes the
threshold-capped field explicitly. A multiplicity cap below two is rejected.

\section{A polynomial-time oracle for actual matching completions}
\label{sec:oracle}

\subsection{Diagram and normalization conventions}
A PD crossing is a cyclic four-tuple of edge labels. Opposite slots are the
same strand. The two smoothings are
\[
 0:(0,1)(2,3),\qquad 1:(0,3)(1,2).
\]
A closed diagram has each edge label exactly twice. The rotation system must
be spherical; arbitrary four-valent incidence data are not a classical link
certificate. The reference constructor checks this using the edge
involution and face permutation. If the crossing shadow has $s$ connected
components and $n$ crossings, the spherical condition is $f=n+2s$, checked
on the resulting orientable combinatorial maps. Empty PD data additionally
specify the number of crossingless components, rather than silently treating
the empty monoidal unit as an unknot.

For a diagram $D$ and state $s\in\{0,1\}^n$, let $c(s)$ count its resolved
circles and $|s|$ its number of one-smoothings. Define the raw reduced bracket
\begin{equation}\label{eq:bracket}
 J_D(q)=\sum_{s\in\{0,1\}^n}(-q)^{|s|}
             (q+q^{-1})^{c(s)-1}.
\end{equation}
For a marked diagram this is its raw reduced quantum Euler polynomial.
Orientation and writhe normalization multiply it by an overall signed
monomial. Such factors must be tracked consistently when summing object
contributions but do not change the norm of a whole completed summand.
The implementation uses~\eqref{eq:bracket} throughout, avoiding objectwise
normalizations that would introduce incompatible phases.

\subsection{The value at one and exponent parity}
\begin{lemma}\label{lem:one}
Orient each component of $D$ arbitrarily. Let $n_-(D)$ be the number of
crossings whose oriented smoothing is state one in the above convention,
and let $\ell(D)$ count link components. Then
\begin{equation}\label{eq:one}
 J_D(1)=(-1)^{n_-(D)}2^{\ell(D)-1}.
\end{equation}
Every exponent of $J_D$ has parity
\begin{equation}\label{eq:parity}
 \epsilon=c(0)-1\pmod2.
\end{equation}
Both quantities can be computed in polynomial time, without enumerating
states.
\end{lemma}
\begin{proof}
At $q=1$, switching a crossing interchanges its zero and one smoothings.
Their signed contributions in~\eqref{eq:bracket} are interchanged with an
overall minus sign. The oriented smoothing bit at that crossing also
changes. Hence $(-1)^{n_-}J_D(1)$ is invariant under crossing switches, in
addition to the usual normalized bracket's Reidemeister invariance.
Crossing switches can make each component descending and separate the
components by an ordering of their overcrossings. The result is an unlink
with the same $\ell$. On an $\ell$-component unlink,
\eqref{eq:bracket} is $(q+q^{-1})^{\ell-1}$, giving~\eqref{eq:one}.
This is the standard $q=1$ specialization, with its raw sign written out.

An exponent contributed by a state is $|s|+c(s)-1-2j$ for some integer $j$.
Changing one smoothing changes $|s|$ by one and changes $c(s)$ by one in
absolute value. Thus $|s|+c(s)-1$ has constant parity, namely that of the
zero state. Edge pairing and opposite-slot pairing compute link components
and a choice of orientations. One smoothing computation gives $c(0)$.
These traversals are polynomial and in fact nearly linear with standard
union--find bookkeeping.
\end{proof}

\subsection{The determinant specialization, including its phase}
Assume for the moment that the crossing shadow is nonempty and connected.
Checkerboard-color its spherical complement, and let $G$ be its black Tait
graph. A crossing gives one edge, including possible loops or parallel
edges. For each edge $e$, let $b_e\in\{0,1\}$ be the smoothing bit that
\emph{does not} join its black regions. Put
\[
 v=|V(G)|,\qquad B=\sum_e b_e,\qquad w_e=(-1)^{b_e}.
\]
Let $L_G$ be the signed weighted Laplacian: a nonloop edge of weight $w$
adds $w$ to both endpoint diagonals and $-w$ to the two off-diagonals. Loops
are omitted because they belong to no spanning tree. A deleted root row and
column are denoted by $L_G^{(r)}$.

\begin{theorem}[Exact determinant continuation value]\label{thm:tait}
For a nonempty connected classical shadow,
\begin{equation}\label{eq:tait}
 J_D(i)=(-i)^{B+v-1}\det L_G^{(r)}.
\end{equation}
If the shadow is disconnected, or there is an additional crossingless
component, then $J_D(i)=0$. For a single crossingless unknot it is one.
\end{theorem}
\begin{proof}
At $q=i$, the circle factor $q+q^{-1}$ is zero. Only states with exactly one
circle contribute to~\eqref{eq:bracket}.

For a subset $T$ of black edges, its planar ribbon neighborhood has
$k(T)$ components and Euler characteristic $v-|T|$. Being planar, its number
of boundary components is
\[
 c(T)=2k(T)-v+|T|.
\]
A spanning forest inequality gives $|T|\ge v-k(T)$. If $c(T)=1$, the displayed
formula forces $1\ge k(T)$, so $k(T)=1$ and $|T|=v-1$. Thus $T$ is a spanning
tree. Conversely a tree has a disk neighborhood and one boundary component.
Connected classical shadows have disk faces, so this ribbon-neighborhood
correspondence applies without hidden holes in a checkerboard region.

For its associated state, an edge outside $T$ contributes smoothing $b_e$,
and an edge inside $T$ contributes $1-b_e$. Therefore
\[
 |s(T)|=B+\sum_{e\in T}(1-2b_e),
\]
and since $|T|=v-1$,
\[
 (-i)^{|s(T)|}=(-i)^{B+v-1}\prod_{e\in T}(-1)^{b_e}.
\]
Sum over spanning trees. The weighted matrix-tree identity gives their
signed weight sum as $\det L_G^{(r)}$.

For completeness, the signed identity follows by writing a grounded
incidence matrix $E$ and $W=\operatorname{diag}(w_e)$. Then
$L_G^{(r)}=EWE^T$. Cauchy--Binet expands its determinant as
$\sum_{|T|=v-1}\det(E_T)^2\prod_{e\in T}w_e$.
An incidence minor is zero unless $T$ is a spanning tree, and is $\pm1$ for
a tree, by leaf elimination. This proof uses no positivity of the weights.

Finally, different shadow components can never be joined by smoothings.
Every state then has at least two circles, and the specialization vanishes.
An extra crossingless component has the same effect. A single circle has
reduced bracket one by definition.
\end{proof}

The unsigned quantity $|J_D(i)|$ is the usual determinant specialization,
but the phase in~\eqref{eq:tait} cannot be discarded when the result is an
object's contribution to a larger formal summand. The implementation uses
pairs of integers for Gaussian integers; no floating-point roots of unity
or numerical determinants occur.

\subsection{Reconstructing the entire residue-four vector}
\begin{corollary}[Two evaluations suffice]\label{cor:fouroracle}
Let $a=J_D(1)$, $\epsilon$ be as in~\eqref{eq:parity}, and
$d=i^{-\epsilon}J_D(i)$. Then $d\in\Z$ and
\begin{equation}\label{eq:reconstruct}
 \sig(J_D)_\epsilon=(a+d)/2,\qquad
 \sig(J_D)_{\epsilon+2}=(a-d)/2,
\end{equation}
with the other two coordinates zero. The two displayed halves are integers.
Thus the residue-four vector of a completed matching can be obtained in
polynomial bit time.
\end{corollary}
\begin{proof}
By parity, only the two indicated residues occur. Let their integer
coefficients be $u,v$. Evaluation at one gives $a=u+v$, and evaluation at $i$
gives $J_D(i)=i^\epsilon(u-v)$. Solving yields~\eqref{eq:reconstruct} and all
integrality statements. The two evaluations were computed polynomially in
Lemma~\ref{lem:one} and Theorem~\ref{thm:tait}; exact integer determinants
have polynomial bit complexity.
\end{proof}

\paragraph{Completion of a matching.}
The suffix PD has boundary edge labels occurring once and internal labels
occurring twice. A matching pairs every boundary label exactly once. Join
paired labels by union--find, relabel the suffix crossings, and validate the
resulting closed PD rotation system. This constructs the actual completion
at the level needed by the state sum. The observer caches the answer for
each distinct matching. Nonclassical completions or inconsistent boundaries
are rejected rather than assigned an arbitrary determinant.

For a component with recovered shifts $q_v$ and homological degrees $h_v$,
the completed vector is
\begin{equation}\label{eq:objectsum}
 \sig(F_i^\bullet T)=\sum_v(-1)^{h_v}q^{q_v}
                     \sig(J_{D_i(a_v)})\pmod{q^4-1}.
\end{equation}
No differentials of the completed suffix are required to evaluate this
Euler expression. Their existence and grading are what justify its use as a
homological lower bound.

\section{Singular-safe boundary compression for signed determinants}
\label{sec:kernel}

\subsection{Terminal identifications and the precise algebraic problem}
Let $L$ be a symmetric integer matrix indexed by $N$ vertices. In the graph
application it is a signed weighted Laplacian. Choose an ordered nonempty
terminal set $\mathcal B$ of size $b$, and let $\mathcal I$ be the remaining
vertices. A partition $\pi$ of the terminals identifies terminals in each
block. Ground the block containing the first terminal.

If $\pi$ has $p$ blocks, put $q=p-1$. Let $Q_\pi$ be the $b\times q$ matrix
whose columns indicate the nongrounded blocks. Terminal rows in the grounded
block are zero. The grounded quotient matrix is exactly
\begin{equation}\label{eq:quotient}
 M_\pi=
 \begin{pmatrix}I&0\\0&Q_\pi\end{pmatrix}^{\!T}
 L
 \begin{pmatrix}I&0\\0&Q_\pi\end{pmatrix},
\end{equation}
where the identity is indexed by $\mathcal I$. For a Laplacian this sums
entries over identified vertices, with root rows and columns removed.
Loops created by identification cancel in the resulting Laplacian.

Writing $A=L_{\mathcal I\mathcal I}$, the usual formula
$\det M_\pi=\det A\det(Q_\pi^T S Q_\pi)$ is legitimate only if $A$ is
invertible. Signed graph interiors can fail that hypothesis even when a
quotient determinant is nonzero. A Moore--Penrose inverse inserted without
an additional border does not repair the determinant identity.

\subsection{A maximal principal block}
\begin{lemma}[Symmetric rank-revealing elimination]\label{lem:principal}
Every symmetric rational matrix $A$ has a nonsingular principal submatrix
$P$ of order $\rk A$. After permuting coordinates,
\[
 A=\begin{pmatrix}P&U\\U^T&V\end{pmatrix},
 \qquad V-U^TP^{-1}U=0.
\]
Such a block can be found by exact symmetric pivots of orders one and two
using $O(N^3)$ rational arithmetic operations for an $N\times N$ matrix.
The empty block is allowed and has determinant one.
\end{lemma}
\begin{proof}
If a diagonal entry is nonzero, use it as a one-dimensional pivot. If all
diagonal entries vanish but some off-diagonal entry $a_{ij}$ is nonzero,
the principal block on $i,j$ is
$\left(\begin{smallmatrix}0&a_{ij}\\a_{ij}&0\end{smallmatrix}\right)$,
with determinant $-a_{ij}^2\ne0$. Symmetric Schur elimination preserves
symmetry on the remaining coordinates. Continue until the residual matrix
is zero. At each pivot the rank drops by precisely the pivot order, so the
selected principal block has order $\rk A$ and is nonsingular. Its residual
Schur complement is zero. Each step updates a quadratic number of entries;
there are at most $N$ steps.
\end{proof}

Apply this lemma only to the interior block, while updating its couplings
to terminals. Let $r=\corank A$. In the order selected interior, residual
interior, terminals, write
\begin{equation}\label{eq:fullblocks}
 L=\begin{pmatrix}
 P&U&C_1\\
 U^T&V&C_2\\
 C_1^T&C_2^T&D
 \end{pmatrix}.
\end{equation}
Define
\begin{equation}\label{eq:RS}
 R=C_2-U^TP^{-1}C_1,\qquad
 S=D-C_1^TP^{-1}C_1,\qquad \alpha=\det P.
\end{equation}
The $r\times b$ matrix $R$ is the essential nullspace coupling. Dropping it
is exactly the error avoided by the next theorem.

\begin{theorem}[Singular-safe terminal kernel]\label{thm:kernel}
For every terminal partition $\pi$,
\begin{equation}\label{eq:kernel}
 \det M_\pi=\alpha\det
 \begin{pmatrix}
 0_{r\times r}&R Q_\pi\\
 Q_\pi^T R^T&Q_\pi^T S Q_\pi
 \end{pmatrix}.
\end{equation}
If $r>p-1$, or if $RQ_\pi$ has rank less than $r$, then this determinant is
zero. In particular $r>b-1$ forces every terminal-partition query to vanish.
Otherwise every potentially nonzero query in~\eqref{eq:kernel} has order at
most $2b-2$.
\end{theorem}
\begin{proof}
Eliminate $P$ from~\eqref{eq:fullblocks} by multiplying on the left and right
by mutually transposed block-unitriangular matrices. Their determinants are
one. The result is
\[
 \begin{pmatrix}P&0&0\\0&0&R\\0&R^T&S\end{pmatrix},
\]
because the interior residual $V-U^TP^{-1}U$ vanishes. Terminal
identification and grounding leave the selected interior coordinates
untouched. The same elimination is therefore valid in~\eqref{eq:quotient},
with the terminal factors replaced by $Q_\pi$. Taking determinants gives
\eqref{eq:kernel}.

If $RQ_\pi$ has rank less than $r$, choose a nonzero vector $z$ in the kernel
of $Q_\pi^TR^T$. Then $(z,0)$ is a kernel vector of the bordered matrix.
This includes $r>q=p-1$. A nonzero query must therefore have $r\le p-1\le b-1$,
so its order $r+p-1$ is at most $2b-2$.
\end{proof}

The theorem is algebraic: it holds for any symmetric integer input, not
only positive Laplacians or planar graphs. A geometric extraction theorem
is a separate matter, treated in Section~\ref{sec:integration}.

\subsection{The critical square law}
\begin{theorem}[Critical terminal partitions have a fixed square class]
\label{thm:square}
In the setting of Theorem~\ref{thm:kernel}, suppose $\pi$ has exactly $r+1$
blocks. Then
\begin{equation}\label{eq:square}
 \det M_\pi=(-1)^r\alpha\,\det(RQ_\pi)^2.
\end{equation}
Consequently all nonzero critical-partition determinants have the same sign
and the same class in $\Q^\times/(\Q^\times)^2$. For integer $L$, the product
of any two nonzero such determinants is a positive integer square.
\end{theorem}
\begin{proof}
Now $Q_\pi$ has $r$ columns. Write $C=RQ_\pi$ and $E=Q_\pi^TSQ_\pi$.
Interchanging the two groups of $r$ columns in
$\left(\begin{smallmatrix}0&C\\C^T&E\end{smallmatrix}\right)$ has sign
$(-1)^{r^2}=(-1)^r$ and yields the block-triangular matrix
$\left(\begin{smallmatrix}C&0\\E&C^T\end{smallmatrix}\right)$.
Its determinant is $\det(C)^2$, including when $C$ is singular.
Equation~\eqref{eq:kernel} proves~\eqref{eq:square}.

Nonzero squares in $\Q$ are positive. Thus the common sign and square class
are those of $(-1)^r\alpha$. For two partitions the product is the square
of the rational number $\alpha\det(RQ_\pi)\det(RQ_\rho)$. It is also an
integer. A rational number whose square is an integer is itself an integer:
write it in lowest terms and compare prime factors of the denominator.
This proves the final assertion.
\end{proof}

The response matrix $S$ does not affect critical queries at all. This gives
an exact shortcut, implemented in \code{TerminalKernel.query}, and a useful
metamorphic test: changing $S$ while preserving the other hypotheses cannot
change a critical answer. The statement concerns graph cofactors. Converting
it to a claim about a family of tangle closures still requires a correct
common-graph extraction and the phase bookkeeping of~\eqref{eq:tait}.

\subsection{Sharpness and a minimal counterexample}
\begin{example}[A singular interior with nonzero determinant]\label{ex:singular}
Take one interior vertex $u$ and terminals $t_0,t_1$, with edge weights $+1$
on $ut_0$ and $-1$ on $ut_1$. In this order the Laplacian is
\[
 L=\begin{pmatrix}0&-1&1\\-1&1&0\\1&0&-1\end{pmatrix}.
\]
The interior block is zero. For distinct terminals, grounding $t_0$ leaves
$\left(\begin{smallmatrix}0&1\\1&-1\end{smallmatrix}\right)$, of determinant
$-1$. Identifying both terminals gives determinant zero. Thus singularity of
$A$ does not make all completions vanish. This is already a planar signed
graph, not an artifact restricted to nonplanar test matrices.
\end{example}

\begin{proposition}[The corank gate is sharp]\label{prop:sharp}
For every $b\ge2$ there is a planar signed graph with $b$ terminals and
interior corank $b-1$ whose distinct-terminal determinant is nonzero.
Hence the threshold $r>b-1$ cannot be improved using only $r$ and $b$.
\end{proposition}
\begin{proof}
Put $r=b-1$. Use terminals $t_0,\ldots,t_r$ and interior vertices
$u_1,\ldots,u_r$. Add edges $u_jt_0$ of weight $+1$ and $u_jt_j$ of weight
$-1$. The graph is a subdivided star and hence planar. Each interior diagonal
is zero and there are no interior off-diagonals, so $A=0_r$ has corank $r$.
For the distinct-terminal partition, $RQ=I_r$ up to a simultaneous sign
convention, and~\eqref{eq:square} gives $(-1)^r$. Equivalently, the graph
itself is its unique spanning tree and has that signed product.
\end{proof}

This proves sharpness of the \emph{vanishing threshold}, not a lower bound
on every possible representation of the determinant function. In particular
one must not infer an information-theoretic necessity for a $2b-2$ matrix
from the size of the displayed bordered formula.

\subsection{Batch arithmetic cost and exact bit sizes}
\begin{corollary}[Batch cost]\label{cor:batch}
For a fixed symmetric integer matrix of order $N$, $b$ terminals, and $M$
terminal-partition queries, the kernel can be prepared and queried in
\[
 O(N^3+Mb^3)
\]
rational arithmetic operations. After the all-zero corank test, its numeric
query core needs $O(b^2)$ rational entries. Preprocessing may use $O(N^2)$
workspace; certificate and original-input storage are separate.
\end{corollary}
\begin{proof}
Lemma~\ref{lem:principal} gives cubic preprocessing. If $r>b-1$, every answer
is zero and the numeric query data can be replaced by that certified flag.
Otherwise $R$ and $S$ have $O(b^2)$ entries. Since $Q$ is a partition
indicator, aggregation into $RQ$ and $Q^TSQ$ costs $O(b^2)$ additions. The
remaining determinant has order at most $2b-2$ and costs $O(b^3)$ rational
operations. Reading each partition costs $O(b)$ and is absorbed.
\end{proof}

To turn this into a bit statement, suppose $|L_{ij}|\le H$. Hadamard's
inequality bounds a minor of order $t\le N$ by $(\sqrt t\,H)^t$. Its bit
length is $O(N(\log(N+2)+\log(H+2)))$. The entries in $R,S$ are ratios of
bordered minors to $\det P$; their reduced numerators and denominators
therefore have polynomially bounded bit length. Exact query elimination
involves minors of matrices of order $O(b)$ with those rational entries,
so its bit cost is polynomial in $N,b,\log(H+2)$ per query. No arithmetic
operation on a long integer is assumed to cost one bit operation.

For signed unit-weight Tait graphs with at most $n$ edges, one may take
$H\le n$ and $N=O(n)$, giving polynomial bit cost in the original crossing
count. The reference code uses Python integers and \code{Fraction}, an
integer Bareiss determinant for the direct oracle, and exact symmetric
Schur elimination for the kernel. It retains extra index and block data for
independent certificate checking, including in all-zero cases; thus the
certificate-retaining object is not advertised as a minimal $O(b^2)$-entry
serialization in every degenerate case.

\section{Integration contracts and executable algorithms}
\label{sec:integration}

\subsection{The general read-only observer}
The package provides an observer for the stable in-memory shape of
\code{FastScan} and \code{ComponentScan}. It reads matching identifiers,
homological degrees, nonzero differential rows, and the algebra's matching
pairs. When available, it also reads owner identifiers and multiplicity
weights. It neither changes the differential nor imports the production
scanner to run its local unit tests.

\begin{lstlisting}
observe(valid_relative_scan, actual_nonempty_suffix, retained_mark):
    check that the mark remains in the suffix
    find connected components of the live nonzero differential graph
    for each stored component T:
        check homogeneous coefficients and recover relative q-potentials
        for each object a in T:
            complete a with the actual suffix
            validate its closed classical PD
            cache its residue-four vector using q=1 and q=i
        add the signed, q-shifted object vectors within T
        take the integral l1 norm
        multiply by T's nonnegative multiplicity
    return the sum as a reduced-rank lower bound
\end{lstlisting}

The implementation is \path{detshadow/continuation.py}. It explicitly rejects
an absent mark, mismatched boundaries, a nonhomogeneous coefficient, an
inconsistent quantum potential, and invalid differential endpoints. These
checks are necessary but not a complete proof-carrying reconstruction of
$C_i$ from the original knot. An arbitrary object with scanner-like fields
is not automatically a knot certificate.

The returned record distinguishes the residue-four, Euler, and determinant
character lower bounds. It always reports that it does not certify an
unknot. Calling code may reject a \emph{validated knot} when the first bound
is greater than one, provided the relative provenance contract has been
met. A resource limit is not a topological result.

\subsection{When one terminal kernel may be reused}
For a batch of completed matchings, one would like a fixed signed graph
$G_i$ with terminals, and for each matching $a$ a partition $\pi_a$ and phase
$\eta_a$ such that
\begin{equation}\label{eq:geometrycontract}
 J_{D_i(a)}(i)=\eta_a\det L_{G_i/\pi_a}^{(\mathrm{root})}.
\end{equation}
A certificate for this must justify the graph correspondence and phase,
not merely exhibit a numerically convenient partition. A common-disk tangle
picture often motivates terminal identifications of black regions, but an
arbitrary PD frontier does not automatically supply that picture. Closed
regions, disconnected shadows, nesting, and the retained marked arc must
all be handled explicitly.

This package supplies a conservative \emph{verifier} for a proposed quotient
identification. It independently constructs the actual completion's Tait
cofactor and checks its equality, under a supplied row-and-column
permutation, with the proposed quotient cofactor. The raw phase is obtained
from the actual PD, not guessed from an unsigned determinant. Combined with
the separate kernel block-identity verifier, this certifies the numerical
bridge. It does not automatically \emph{find} a useful fixed graph or the
terminal partitions.

That verification has $O(N^2)$ matrix-comparison overhead per query, in
addition to constructing the actual Tait graph. Therefore the pure
$O(N^3+Mb^3)$ batch theorem is not advertised as the cost of an automatically
extracted arbitrary-PD observer. With per-query cofactor certificates the
additional $O(MN^2)$ work must be charged. A proved boundary-only extractor
would be needed to remove that term. The general observer in the package
uses fresh actual completions and remains valid without any common-graph
hypothesis.

\subsection{Certificate layers}
There are three logically distinct certificates:
\begin{description}[style=nextline]
\item[Relative chain provenance.]
The original knot diagram, scan order, valid local operations, cancellation
maps, and direct-sum ancestry justify the processed complex. The present
observer takes this as an input contract of the existing scanner.
\item[Completion geometry.]
A matching covers the actual suffix boundary, the completed rotation system
is classical, and a retained mark makes reduced completion legitimate.
For batch reuse, an additional quotient identification and phase are needed.
\item[Exact linear algebra.]
The selected principal block is nonsingular, the residual interior Schur
block is zero, and $R,S,\alpha$ satisfy~\eqref{eq:RS}. The function
\code{verify\_kernel} checks these identities using a separate block solve,
not by trusting the builder's stored pivot history.
\end{description}
Neither a checked determinant nor a checked $d^2=0$ identity substitutes for
the first two layers. Conversely, once those layers are valid, the observer
does not need to enumerate the suffix crossing cube.

\subsection{Diagnostic runner and release boundaries}
The supplied \path{integration/upstream_probe.py} accepts a local
\path{fast/} checkout. It observes proper prefixes of natural and shuffled
orders, then completes each scanner and compares all lower bounds with its
reduced rank. It also compares small cases with the independent cube,
checks $d^2=0$, and records imported source hashes. Object exhaustion yields
an explicit \code{UNKNOWN} record.

This diagnostic was \emph{not run against an upstream checkout in this
session}. It is not a production performance patch. In particular, the
reference determinant observer is not yet interruptible inside a long
integer elimination, so it must not be inserted into a hard-deadline
production path without a budget policy or process boundary. The safe
initial integration is opt-in, read-only observation with regression checks,
not a new default recognizer.

\section{Stopping-prefix complexity and the quasi-polynomial target}
\label{sec:complexity}

\subsection{Execution parameters rather than an unproved structural promise}
Let $n$ be the original crossing count, and fix a scan order. Let $j$ be a
proper prefix at which a marked residue-four witness is found. Define $W$
to be the largest frontier size among the \emph{visited} scan stages, and
let $R\ge1$ bound the physically allocated object count at every such stage,
including expansion before reduction, cancelled slots still stored in
arrays, and representatives before sharing. Coefficient support sizes and
exact multiplicity bit lengths are charged separately below.

These are actual execution parameters. Replacing $R$ by the number of
matching types, or by a minimal survivor count that omits temporary
expansion, would not justify the same bound. Likewise, existence of a good
order does not give an algorithm to find one; order generation must be
supplied or included in the cost.

\begin{theorem}[Controlled successful-prefix bound]\label{thm:complexity}
Assume the existing exact scanner is run by explicit crossing extension and
ordinary finite scalar-unit elimination, with polynomial bookkeeping in its
physical object count and coefficient operations costing $2^{O(W)}$ times
polynomial factors. Suppose its stagewise decompositions satisfy the marked
and grading hypotheses proved above. Using the general actual-completion
observer at at most $n$ visited stages adds no dependence on the maximum
frontier of an unvisited suffix, and the total work through a successful
prefix is bounded by
\begin{equation}\label{eq:cost}
 T_{\mathrm{stop}}\le \poly(n,R)\,2^{O(W)}
\end{equation}
in bit operations, with an additional separately charged order-construction
cost. In particular, if $W=O(\log^2(n+2))$ and
$\log R=O(\log^2(n+2))$, the stopping-prefix work is
$2^{O(\log^2(n+2))}=n^{O(\log n)}$.
\end{theorem}
\begin{proof}
At each visited stage there are at most $R$ objects and $O(R^2)$ potential
differential entries. At a frontier with $W$ endpoints, an overlay has at
most $W/2$ circles, so a fully stored binary coefficient has at most
$2^{W/2}$ bits. The assumed finite elimination and its bookkeeping are
therefore bounded by a polynomial in $n,R$ times $2^{O(W)}$. This is a
conservative bound, not a prediction of sparse performance.

Finding differential components and recovering shifts requires traversal of
these explicit entries and, for validation, their coefficient supports.
The same factor covers it. A matching completion retains at most $n$
crossings and is built from polynomial-size PD data. Its connectivity,
checkerboard graph, signed integer cofactor, and residue reconstruction all
cost polynomial bit time in $n$. There are at most $R$ distinct matching
queries per observation and at most $n$ observations.

Each crossing has two smoothings and can create at most two closed circles
in a local gluing. Delooping therefore creates at most eight children per
old object. The virtual expanded object count through $n$ crossings is at
most $8^n$ starting from the unit. Exact multiplicities consequently have
$O(n)$ bits. An object-completion Euler value has at most exponential
magnitude in $n$, and adding at most $R$ such values uses
$O(n+\log(R+1))$ bits. Recovering shifts and applying them modulo four also
uses polynomial-size integers. Thus neither signed accumulation nor
multiplicity arithmetic hides another uncharged exponential bit cost.

The suffix is evaluated by a determinant of polynomial order, not by
scanning its crossing cube. Hence its possible scan-frontier width does
not enter this observer cost. Combining these estimates proves
\eqref{eq:cost}. Substitution of the two logarithmic-square hypotheses
proves the final statement.
\end{proof}

The inherited Euler observer already has a direct connectivity route. The
new point is not that every old suffix-width dependence is newly removed;
it is that the \emph{stronger marked determinant/residue observer} also has
a polynomial completion oracle, while its erased grading can be recovered
from the existing keys.

\subsection{What this theorem does not establish}
Equation~\eqref{eq:cost} does not force $B_4$ to become larger than one before
$R$ or $W$ grows. It also does not show that a nontrivial knot must be exposed
by this observer while the mark remains in a proper suffix. An unknot never
has such a rejection witness. It still needs a separate complete certificate
or the existing exact backend. The current source and this package provide
no universal bound placing all inputs in the stated quasi-polynomial regime.

A final reduced complex with zero differential can be split into individual
graded generators, at which point its sum of shadow norms is exactly its
rank. This mathematical observation is not a shortcut: obtaining that
complex may already have incurred exponential work. It also differs from
illegally dividing arbitrary unmarked closed summands, as warned in
Remark~\ref{rem:divide}.

For completeness, full reduced rank one over $\F$ still suffices for
unknot recognition. Universal coefficients bound the reduced rational rank
by the $\F$ dimension. The rational rank is at least one, for instance by
its normalized Euler value at one. Thus $\F$ rank one forces rational rank
one, and Kronheimer--Mrowka apply~\cite{km}. None of these implications says
that a partial lower bound of one is sufficient.

\subsection{A reasonable performance policy}
The observer should initially be used where the maintained algorithm has
already discovered several distinct summands or where cached completions
make its cost modest. An optional policy can skip a probe after its budget
is exhausted and continue the original scanner. Such skipping loses a
possible early certificate but cannot change the represented complex.
The lower bound from an earlier successful probe remains valid.

To evaluate actual improvement, measure the added determinant cost, the
first stage exposing $B_4>1$, and the remaining scan work avoided. Reporting
only a faster determinant kernel would omit the decisive end-to-end
question. Conversely, large virtual multiplicities can make a small
positive template contribution decisive, but only after that contribution
has been justified by marked completion.

\section{Validation, measurements, and reproducible artifacts}
\label{sec:validation}

\subsection{Independent checks actually executed}
All computations in the delivered local package use the Python standard
library. The independent cube is an exponential validation oracle and is
not called by the fast matching-specialization path. The completed checks
are summarized in Table~\ref{tab:checks}; full inputs and outcomes are
retained in the JSON records.

\begin{table}[ht]
\centering
\begin{tabular}{lr}
\toprule
Check & Completed count\\
\midrule
Unit-test methods & 35\\
Diagram shadow comparisons & 701\\
Explicit smoothing states in those comparisons & 133,923\\
Reduced homology comparisons & 83\\
Columns checked for $d^2=0$ & 62,904\\
Signed random graphs & 160\\
Graphs with singular interior blocks & 57\\
Terminal-quotient determinant comparisons & 4,000\\
\bottomrule
\end{tabular}
\caption{Executed local checks. Counts refer to the delivered run, not to
an upstream production test suite.}\label{tab:checks}
\end{table}

The diagram experiment uses seed 20261007 and 700 generated braid closures
on two through six strands with at most ten crossings, including links,
unused strands, and disconnected shadows. It adds the repository's
11-crossing Conway PD~\cite{conwayinput}. Every fast residue vector is compared with the full
state sum, which independently expands the factors
$(q+q^{-1})^{c(s)-1}$. A bounded subset is compared with a separately built
marked reduced differential and exact binary homology. In that oracle the
marked circle is forced to carry $x$, with its conventional reduced grading
shift. Merge and split maps are implemented directly from the Frobenius
rules; every generated two-step differential column is checked.

The Conway input has reduced $\F$ rank 33 in the independent computation,
while its global residue-four norm is one. The positive three-braid
$(\sigma_1\sigma_2)^5$ has reduced rank seven and again global residue norm
one. These are controls demonstrating the incompleteness of the global
specialization. They do \emph{not} demonstrate an early component-stage
speedup on these knots, because no upstream scan of them with the observer
was executed here.

The graph experiment compares every kernel query with an independently
constructed quotient matrix and an integer Bareiss determinant. It includes
singular interiors and checks the kernel by an independent block-identity
solve. Unit tests additionally cover zero-dimensional determinants,
two-dimensional pivots with zero diagonals, all partitions of small terminal
sets, the sharp-corank family through $r=12$, the critical square law,
coefficient homogeneity, potential cycle failures, and deliberately corrupted
algebraic or geometry certificates.

Finite computations support the implementation and expose normalization
errors; they are not substitutes for the written proofs. Passing these
checks is not a Lean formalization, an exhaustive theorem-priority review,
or a complete topological certificate for arbitrary externally supplied
scanner objects.

\subsection{Paired exact determinant-kernel benchmark}
The benchmark compares two fresh direct controls, A and B, with a compressed
arm. Each direct arm rebuilds every quotient matrix and computes its exact
Bareiss determinant. The compressed arm includes both kernel preprocessing
and all its queries. Graph construction precedes the timed arms; independent
certificate verification is excluded from all arms. Arm order is randomized
in each of three paired rounds. Every integer answer is checked for equality.
Each case has six terminals and 96 queries.

\begin{table}[ht]
\centering\small
\begin{tabular}{rrrrrr}
\toprule
$N$ & $r$ & Direct A (s) & Compressed (s) & A/compressed & A/B\\
\midrule
32 & 0 & 0.1054 & 0.0189 & 5.52 & 1.002 \\
56 & 0 & 0.5727 & 0.0830 & 6.90 & 1.011 \\
80 & 0 & 1.8632 & 0.2928 & 6.43 & 1.031 \\
56 & 1 & 0.6086 & 0.0864 & 7.05 & 0.971 \\
\bottomrule
\end{tabular}
\caption{Exact determinant-kernel measurements. Times are medians; ratios are
medians of paired ratios, not ratios of medians. The compressed arm includes
preprocessing. Each case uses six terminals and 96 queries.}\label{tab:bench}
\end{table}

The measurements are from CPython 3.13.5 on the recorded Linux x86-64 host.
They are small-sample, host-dependent kernel measurements, not fitted
asymptotic exponents or hardware-independent guarantees. The A/A ratios
measure ordinary timing variability; three rounds are not a confidence
interval. The signed graph family is not asserted to be the frontier graph
of a measured knot scan. In particular the table does not establish that
recognizing a knot is faster by its listed factors.

\subsection{Package map and reproduction}
The package contains the self-contained article source and compiled PDF,
six implementation modules, unit tests, complete validation and benchmark
JSON, a worked singular-kernel certificate, source provenance, and integration
notes. The important paths are:
\begin{center}\small
\begin{tabular}{p{0.40\linewidth}p{0.53\linewidth}}
\toprule
Path & Purpose\\
\midrule
\path{detshadow/continuation.py} & Quantum recovery and read-only component observer.\\
\path{detshadow/diagram.py} & PD validation, signed Tait phase, residue oracle.\\
\path{detshadow/linalg.py} & Exact determinants, singular kernel, independent verification.\\
\path{detshadow/geometry_certificate.py} & Proposed quotient-to-actual-PD verifier.\\
\path{detshadow/cube.py} & Independent reduced homology validation oracle.\\
\path{experiments/} & Seeded comparisons and paired kernel benchmark.\\
\path{integration/upstream_probe.py} & Local-checkout diagnostic, supplied but not executed here.\\
\path{results/} & Raw evidence, not merely aggregate speed claims.\\
\bottomrule
\end{tabular}
\end{center}

From the package root:
\begin{lstlisting}
python -m unittest discover -s tests -v
PYTHONPATH=. python experiments/validate.py
PYTHONPATH=. python experiments/benchmark.py
cd paper && sh build.sh
\end{lstlisting}
No knot-library dependency is required. Running the experiments overwrites
the corresponding JSON with fresh local measurements; the delivered PDF
records the delivered run. Preserve those records before producing a
comparison. The article's TeX file embeds its tables and bibliography, so
it can also be compiled independently.

\section{Further research questions with concrete success criteria}
\label{sec:questions}

\subsection{Immediate correctness and integration questions}
\paragraph{1. Complete upstream prefix regression.}
Run the supplied diagnostic on a pinned complete checkout, with natural,
greedy, and shuffled orders. The immediate success criterion is agreement
with full reduced ranks, valid recovered potentials at every tested stage,
and monotonicity whenever the component decomposition evolves by refinement.
Tests should include determinant-one knots and RI/RII-expanded unknots.
A failed potential equation should be diagnosed as a grading or interface
failure, not bypassed.

\paragraph{2. A certified terminal-partition producer.}
Construct a map from a suitable frontier embedding to one signed suffix
graph and the partition and phase for every matching. Prove exactly when
black-region identifications are sufficient and when extra region or
connectivity data are needed. A useful result would give $\poly(b)$
per-matching extraction and a checkable embedding certificate, rather than
only an $O(N^2)$ cofactor comparison. Arbitrary PD frontiers must not be
silently treated as one disk.

\paragraph{3. Interruptible exact arithmetic and inference budgets.}
Add cancellation-safe interruption points or isolate determinant probes in
bounded worker processes. Record preprocessing, rational bit sizes,
coefficient support scans, and observer time separately from scanner time.
The success criterion is that exhausted optional inference leaves the
original complex untouched and never becomes an unknot verdict.

\subsection{Questions about exposing useful witnesses}
\paragraph{4. Early residue witnesses for determinant-one knots.}
On the repository's Conway, Kinoshita--Terasaka, and determinant-one stress
inputs, determine whether literal component splitting exposes $B_4>1$
substantially before the complete backend. A positive result needs paired
end-to-end measurements and the associated stage records. A negative result
would be equally informative if it identifies why summands remain coupled.

\paragraph{5. Root and order optimization.}
The retained root crossing constrains the final part of an order. Can one
choose a root and order jointly to minimize observed cost until a witness
appears? A proved target would be a constructive family theorem, not just
the existence of a good order found by exhaustive search. Certificate
quality and frontier width need not be optimized by the same order.

\paragraph{6. Coherence after local simplification.}
Proposition~\ref{prop:defect} gives an exact target: eliminate overlap between
even and odd homological parity in each quantum residue of completed
summands. Identify a nontrivial class of genuine knot tangles for which
local simplification guarantees this coherence. Such a theorem would turn
the residue observer from a lower bound into an exact rank computation on
that class. The abstract dotted-block example is not such a realization
theorem.

\subsection{Stronger algebraic and structural questions}
\paragraph{7. Dynamic singular kernels.}
As the scan moves, the suffix graph changes. Can $P,R,S$ be updated under
crossing deletion and terminal promotion without rebuilding from scratch,
while certifying rank changes and preserving polynomial bit bounds? An
algorithm assuming constant corank would need a separate mechanism for
rank-drop events; those are precisely where an ordinary inverse update
fails.

\paragraph{8. Topological meaning of the critical square class.}
When a fixed-graph tangle family satisfies the quotient contract, the
critical-partition determinants share one rational square class. Determine
which tangle families realize large corank and whether their square-class
constraint yields useful pruning before exact determinant computation.
Any statement about link determinants must remove raw phases correctly and
handle zero determinants separately.

\paragraph{9. Higher cyclotomic residue observers.}
Replace $q^4-1$ by another explicitly tractable quotient. The algebraic
norm proof extends immediately, but the matching-completion oracle may no
longer be polynomial-time. The substantive target is a stronger observer
with a proved exact evaluation algorithm and bit-cost bound, not simply
more sample points evaluated exponentially.

\paragraph{10. Physical-size bounds and witness exposure.}
Prove, for a meaningful family of input diagrams, that a useful marked
witness is exposed while $W$ and $\log R$ stay polylogarithmic. This is the
missing bridge to a family-level quasi-polynomial recognition theorem.
Bounds on coefficient algebra dimension or matching-type count alone do
not supply it. A strong negative result would construct families forcing
large physical representation before every witness in this observer class.

\paragraph{11. Combining windows, twists, and component observers.}
The new observer can in principle share a valid intermediate complex with
the existing window and twist mechanisms. A composition theorem must track
which guard degrees, marked arcs, and decomposition maps survive the
handoff. A claimed combined speedup should be measured against each existing
mechanism separately as well as against the default pipeline.

\paragraph{12. Formal proof and minimal certificate development.}
The matrix identities, square law, residue norm properties, and graph
potential uniqueness are good initial formalization targets. The harder
step is connecting scanner operations to a marked graded tangle category.
A small numerical certificate is not enough until that semantic bridge is
also trusted or verified.

\section{Conclusion}
The completed continuation supplies a stronger exact observer and a
singular-safe determinant engine. Its most useful interface result is that
homogeneous coefficients recover enough quantum information from existing
ungraded connected-component keys. Its most concrete new algebraic
consequence is the critical-partition square law, together with a sharp
corank gate and a batch determinant formula valid for signed weights.

These results isolate a tractable portion of the recognition problem:
completed summand obstructions can be evaluated without expanding the
suffix cube, and shared graph interiors can be compressed without assuming
invertibility that signed graphs do not possess. The remaining global
question is whether genuine knot computations expose decisive summands
before physical representation growth becomes large. The package provides
proofs, executable checks, and explicit integration boundaries for studying
that question, rather than treating a local kernel improvement as a proof
of unrestricted quasi-polynomial recognition.

\appendix
\section{Worked phase and square-law examples}
\label{app:worked}

\subsection{A trefoil state-sum check}
For the package's positive two-strand word $(\sigma_1)^3$, the raw reduced
polynomial computed by the independent state sum is a signed monomial
normalization of the usual trefoil polynomial. Rather than making an
orientation-dependent table a convention, the test checks the complete
Laurent dictionary against the signed Tait phase and against all four
residues. The reduced rank is three, the absolute Euler value at one is one,
and the determinant character norm is three. Thus the initial residue
observer already detects this knot, although the production pipeline's
existing determinant filter would do so more cheaply. This example validates
normalization; it is not evidence of a new pipeline advantage.

\subsection{A square-class family with a nontrivial scale}
Start with $r$ balanced interior vertices: each is incident only to terminals,
and the sum of its signed incident weights is zero. Thus their interior
block is zero. Add an additional interior leaf attached to the root
terminal by a weight-two edge, or equivalently by two parallel unit edges.
Its selected principal block is $P=(2)$, so $\alpha=2$. For any critical
partition into $r+1$ terminal blocks,
\[
 \det M_\pi=(-1)^r\,2\,\det(RQ_\pi)^2.
\]
In this construction $R$ is integral. The nonzero absolute determinants
are twice integer squares, regardless of arbitrary additional terminal
edges in the response block. This is the family used to test the critical
shortcut for $r=1,2,3$ with five terminals and all critical partitions.
No factorization algorithm or probabilistic square test is involved in the
formula.

\section{Proof and implementation failure checklist}
\label{app:checklist}
\begin{longtable}{p{0.33\linewidth}p{0.60\linewidth}}
\toprule
Tempting shortcut & Required correction\\
\midrule
\endhead
Discard quantum information completely & Recover edge-weight potentials within each connected component; check homogeneity and cycles.\\[3pt]
Use arbitrary unmarked closed summands & Retain a mark outside processed operations before taking reduced completion.\\[3pt]
Use $p_a(-1)$ as multiplicity & Count copies nonnegatively with $p_a(1)$; shifts act only before a shift-invariant norm.\\[3pt]
Saturate signed Euler coordinates & Accumulate exactly; cap only nonnegative contributions after the norm.\\[3pt]
Replace a complex by disconnected vector records & Require genuine relative decomposition and ancestry, not only the appearance of a sparse matrix.\\[3pt]
Drop the phase of $J_D(i)$ & Retain $(-i)^{B+v-1}$ until all terms of a component are summed.\\[3pt]
Use a Tait cofactor on a disconnected shadow & Return the zero specialization, or justify a separately correct disconnected formula.\\[3pt]
Set all queries to zero when $A$ is singular & Keep the residual nullspace coupling $R$; use the bordered formula.\\[3pt]
Use a pseudoinverse as an ordinary inverse & Prove the relevant singular determinant identity; the missing border cannot be omitted.\\[3pt]
Apply a batch graph kernel to every matching & Establish the fixed-graph quotient and phase contract for that geometry.\\[3pt]
Call a parameterized stopping bound universal & Prove constructive bounds on visited physical size and witness exposure for the intended input class.\\[3pt]
Report kernel speedup as recognition speedup & Run the complete opt-in and baseline pipelines with paired measurements.\\
\bottomrule
\end{longtable}

\clearpage
\begin{thebibliography}{99}
\bibitem{khovanov}
M. Khovanov, \emph{A categorification of the Jones polynomial},
1999, \href{https://arxiv.org/abs/math/9908171}{arXiv:math/9908171}.

\bibitem{bn-tangles}
D. Bar--Natan, \emph{Khovanov's homology for tangles and cobordisms},
2004, \href{https://arxiv.org/abs/math/0410495}{arXiv:math/0410495}.

\bibitem{bn-fast}
D. Bar--Natan, \emph{Fast Khovanov Homology Computations},
2006, \href{https://arxiv.org/abs/math/0606318}{arXiv:math/0606318}.

\bibitem{km}
P. B. Kronheimer and T. S. Mrowka,
\emph{Khovanov homology is an unknot-detector},
2010, \href{https://arxiv.org/abs/1005.4346}{arXiv:1005.4346}.

\bibitem{boninger}
J. Boninger, \emph{The Jones Polynomial from a Goeritz Matrix},
2021, \href{https://arxiv.org/abs/2110.03082}{arXiv:2110.03082}.

\bibitem{kron}
F. D\"orfler and F. Bullo,
\emph{Kron Reduction of Graphs with Applications to Electrical Networks},
2011, \href{https://arxiv.org/abs/1102.2950}{arXiv:1102.2950}.

\bibitem{lackenby}
M. Lackenby, \emph{Incompressible surfaces, hierarchies and unknot recognition},
25 July 2026, \href{https://arxiv.org/html/2607.23350v1}{arXiv:2607.23350v1}.

\bibitem{repo}
V. Reshetnikov, \emph{ProveIt: Topology/UnknotRecognition}, source repository,
\href{https://github.com/VladimirReshetnikov/ProveIt/tree/8e5edcedb580523bdd4946233f7d0b42a97901d4/Topology/UnknotRecognition}{observed October 2026 revision}.
See \path{provenance.json} for the exact per-file audit scope.

\bibitem{component}
ProveIt, \emph{component\_scan.py},
\href{https://github.com/VladimirReshetnikov/ProveIt/blob/8e5edcedb580523bdd4946233f7d0b42a97901d4/Topology/UnknotRecognition/fast/fastunknot/component_scan.py}{pinned source}.

\bibitem{grading}
ProveIt, \emph{audit\_grading.py},
\href{https://github.com/VladimirReshetnikov/ProveIt/blob/dfc762327721d9d8040eb62a36a4ab20c005b9e5/Topology/UnknotRecognition/fast/audit_grading.py}{pinned grading diagnostic}.

\bibitem{conwayinput}
ProveIt, \emph{examples/conway.json},
\href{https://github.com/VladimirReshetnikov/ProveIt/blob/8e5edcedb580523bdd4946233f7d0b42a97901d4/Topology/UnknotRecognition/fast/examples/conway.json}{pinned Conway PD}.
\end{thebibliography}
\end{document}
